\documentclass[11pt,a4paper]{article}

\usepackage[T1]{fontenc}
\usepackage[utf8]{inputenc}
\usepackage[margin=1in]{geometry}
\usepackage{enumitem}
\usepackage{titlesec}
\usepackage{xcolor}
\usepackage{fancyhdr}
\usepackage[hidelinks]{hyperref}

\definecolor{accent}{HTML}{B45309}
\definecolor{ink}{HTML}{1F2933}
\definecolor{muted}{HTML}{6B7280}
\definecolor{answerink}{HTML}{374151}

\titleformat{\section}
  {\normalfont\large\bfseries\color{accent}}
  {\thesection.}{0.6em}{}
\titlespacing*{\section}{0pt}{1.8em}{0.9em}

\pagestyle{fancy}
\fancyhf{}
\renewcommand{\headrulewidth}{0.4pt}
\fancyhead[L]{\small\color{muted}PlacementPrep}
\fancyhead[R]{\small\color{muted}Panel Questions with Answers}
\fancyfoot[C]{\small\color{muted}\thepage}

\newlist{questions}{enumerate}{1}
\setlist[questions]{
  label=\textbf{\color{accent}Q\arabic*.},
  leftmargin=*,
  labelsep=0.6em,
  itemsep=1.05em,
  parsep=0pt,
  topsep=0.5em
}

\newcommand{\code}[1]{\texttt{\small #1}}

\newcommand{\ans}[1]{%
  \par\vspace{0.32em}%
  \begingroup
  \small\color{answerink}%
  \leftskip=1.1em
  \noindent\textbf{\color{accent}A.}\hspace{0.45em}#1\par
  \endgroup
  \vspace{0.1em}%
}

\title{\vspace{-2.2em}\textbf{PlacementPrep}\\[0.35em]
  \large\normalfont Panel Questions with Answers --- Tech Stack \& Design Choices}
\author{}
\date{}

\begin{document}
\maketitle
\thispagestyle{fancy}
\vspace{-3em}

\noindent\color{muted}\small
Fifty questions an evaluation panel is likely to raise, each with a model answer.
The emphasis throughout is on \emph{why} a decision was taken and what it cost,
rather than on recalling framework definitions, since that is where most viva
marks are won and lost. Where a decision has a genuine weakness, the answer says
so --- a defended trade-off scores better than a denied one.
\color{ink}\normalsize
\vspace{0.4em}

\section{Architecture and Project Structure}
\begin{questions}

\item The repository is split into \code{frontend/}, \code{backend/} and
\code{database/}. Why is the database a standalone package rather than a module
inside the backend, and what does that separation buy you?
\ans{So the schema has exactly one definition and no dependency on the web
framework. The \code{database} package holds the SQLModel tables and the engine,
reads its connection string from the environment, and imports nothing from
FastAPI --- meaning a migration script, an analytics job or a second service can
use the same tables without dragging in the API. The cost is that the Docker
build context must be the repository root, and the backend has to bootstrap
\code{sys.path} to locate the package.}

\item The Docker build context is the entire repository rather than just
\code{backend/}. Explain why, and what the trade-off is in image size and build
time.
\ans{The image needs both \code{backend/} and \code{database/}, and Docker cannot
reference files above its build context, so the context has to be the common
parent. The trade-off is that the whole tree --- including the iOS app, which the
server has no use for --- is sent to the daemon on every build. A
\code{.dockerignore} at the root excludes it, so the transferred context stays
small; without that file, build times would degrade noticeably.}

\item The Xcode project file is generated by XcodeGen and is deliberately
excluded from version control. What problem does that solve on a team, and what
new failure mode does it introduce?
\ans{The \code{.xcodeproj} is a large machine-generated file that conflicts on
almost every parallel branch and is effectively unmergeable. Generating it from a
short \code{project.yml}, with sources picked up by directory tree, means adding a
file changes no project configuration at all. The new failure mode is that
\code{xcodegen generate} must be re-run after any file is added, deleted or
renamed --- otherwise Xcode simply does not see it, and the error looks like a
missing symbol rather than a stale project.}

\item The iOS application ships with zero third-party dependencies. Was that a
constraint or a choice? What did you have to build yourself as a result, and when
would you abandon that position?
\ans{A choice. It meant hand-writing an RFC~4180 CSV reader, the networking and
Keychain layers, an SVG path parser for the brand marks, and the entire design
system. In return there is no dependency resolution step, no supply-chain
exposure, no third-party breakage on an OS upgrade, and a smaller binary. I would
abandon it for anything security-critical that I am not qualified to write ---
cryptography above the Keychain API being the obvious case --- or for a charting
library, where the work is large and the risk is low.}

\item Quiz and company content is bundled into the app and parsed on the device,
while authentication and interviews go over HTTP. What rule did you use to decide
which side of that line a feature belongs on?
\ans{Anything static, identical for every user, and needed offline is bundled;
anything per-user, secret-bearing, or requiring a language model is server-side.
Quiz JSON and company lists never vary by user and must work on a train with no
signal, so shipping them costs a few megabytes and removes a whole class of
failure. Authentication needs a signing secret and interviews need a model API
key, and neither can live in a binary a student can unzip.}

\item If you were asked to add an Android client tomorrow, which parts of this
codebase would be reusable and which would have to be rewritten?
\ans{The entire backend, the database package, the bundled content and the prompt
files carry over untouched --- the API is a plain JSON contract with no iOS
assumptions in it. Everything in \code{frontend/} would be rewritten: the views,
the design system, the observable stores, the audio capture and the on-device
resume extraction, which uses PDFKit and Vision and would become PdfBox and ML
Kit. Roughly speaking the two thirds of the line count that is Swift is lost and
the one third that is Python is kept.}

\end{questions}

\section{iOS Application and SwiftUI}
\begin{questions}[resume]

\item The data stores use the \code{@Observable} macro while the authentication
layer uses \code{ObservableObject} with \code{@EnvironmentObject}. Why are both
present in one codebase, and would you unify them?
\ans{\code{@Observable} is the newer iOS~17 mechanism and gives finer-grained
invalidation: SwiftUI records which properties a view actually read during
\code{body}, so touching one property re-renders only the views that read it.
\code{ObservableObject} invalidates every observer whenever any published
property changes. The authentication layer predates the migration, and yes, I
would unify on \code{@Observable} --- it is a mechanical change and the mixed
idiom is a genuine inconsistency rather than a considered one.}

\item The per-question quiz timer is a cancellable \code{Task} rather than a
\code{Timer}. What compiler-level constraint forced that decision?
\ans{The session model is \code{@MainActor} isolated, but \code{deinit} is not
isolated, so it cannot synchronously reach main-actor state to call
\code{invalidate()} on a stored \code{Timer}. That is a compile error, not a style
preference. A \code{Task} sidesteps it: it is stored as a cancellable handle,
captures \code{self} weakly, and exits on its own when the model is deallocated.}

\item Explain the actor model in this app. Why are the session models marked
\code{@MainActor}, and where does work deliberately run off the main actor?
\ans{The session models and every data store are \code{@MainActor} because they
exist to drive UI state, and annotating the type removes the need to hop actors
on every mutation. Two things deliberately leave the main actor: resume text
extraction runs in a detached task, because Vision optical character recognition
on a multi-page PDF is slow enough to drop frames; and company CSV parsing runs
off the main actor on first open. Both hand back a finished value, so no shared
mutable state crosses the boundary.}

\item The three resource folders are declared as folder references rather than
groups in \code{project.yml}. What breaks if they are declared as groups?
\ans{Folder references land in the built bundle as real subdirectories, which is
what lets the app enumerate them at runtime with
\code{bundle.urls(forResourcesWithExtension:subdirectory:)}. Declared as groups,
the files are flattened into the bundle root, that subdirectory lookup returns
nothing, and the app launches with an empty quiz list and no company data --- and
it fails silently, because an empty directory listing is not an error.}

\item The quiz session is presented as a \code{fullScreenCover} rather than pushed
onto a navigation stack. Justify that choice in terms of user expectations and
state ownership.
\ans{A quiz in progress owns the screen: there should be no navigation bar
offering a back button and no interactive swipe that could dismiss a run
mid-question. A cover gives exactly that, with a single explicit close control.
It also makes state ownership clean --- the session model is created with the
cover and dies with it, so there is no possibility of a stale half-finished
attempt sitting in a navigation stack.}

\item \code{DashboardView} owns the selected tab and passes a \code{Binding} down
into \code{HomeView}. What would go wrong if Home pushed its own screens instead?
\ans{You would end up with two live copies of the same screen. Tapping a Home
shortcut would push a second quiz list onto Home's own stack while the Quiz tab
kept its own, and the two would drift apart in scroll position and filter state.
The tab bar would also lie, showing Home selected while a quiz screen is on
display. Passing a binding upward makes the shortcut switch tabs instead, so
there is one instance of each screen.}

\item Stores are injected via \code{.environment(...)} and inherited by sheets and
covers. How does that affect the SwiftUI preview system, and how did you keep
previews from touching real user data?
\ans{Inheritance means a presented screen does not re-inject anything, which is
convenient in the app but dangerous in previews --- a preview that reached the
real store would read and write the user's actual \code{UserDefaults}. Each store
therefore exposes a preview factory that builds a throwaway instance, several
backed by a \code{UserDefaults} suite named with a fresh UUID. Previews construct
their own instances, so they are isolated by construction rather than by
convention.}

\item The JSON decoder uses a lenient date-parsing strategy rather than the
built-in ISO~8601 strategy. What specific input made the standard strategy fail?
\ans{The backend emits two shapes: timezone-aware ISO~8601 strings, and naive
microsecond-precision strings for datetimes that have been round-tripped through
SQLite, which does not carry a timezone. The stock strategy rejects both --- it
will not accept fractional seconds, and it will not accept a missing timezone.
The lenient strategy tries the plausible formats in turn, which is the pragmatic
fix; the more correct fix is to normalise what the server emits.}

\end{questions}

\section{Design System and User Interface}
\begin{questions}[resume]

\item Colour tokens resolve through an observable theme store rather than being
static literals. Explain the mechanism that lets the whole app repaint on a theme
change without any call site being modified.
\ans{Each token such as \code{Color.ppAccent} is a computed property that reads
the current palette from the observable theme store. Because that read happens
inside a view's \code{body}, SwiftUI's observation machinery registers a
dependency on the store at exactly that point. Changing the theme therefore
invalidates precisely the views that drew a themed colour, and no screen contains
any code that knows a theme exists.}

\item The theme store is a singleton, whereas every other store is injected
through the environment. Why the inconsistency, and is it defensible?
\ans{It is forced by the token design. \code{Color.ppAccent} is a static computed
property, not a view, so it has no environment to read from --- a static context
can only reach a global. Given that the alternative is threading a theme
parameter through every colour reference in the app, the singleton is defensible.
The honest cost is that it is not injectable, so previews and tests all share one
instance and cannot exercise themes in isolation.}

\item There is a separate \code{ppOnAccent} token for marks drawn on top of an
accent fill. Why can that not simply reuse the background token?
\ans{Because it is the theme's ink, and ink is not the same colour on every
theme. On the two dark themes the ink happens to equal the ground colour, so the
two look interchangeable; on the light themes the ink is the text colour and the
ground is paper. Reusing the ground token would paint white marks on a light
amber fill and make every button label and tick invisible.}

\item The Google sign-in mark hardcodes its four brand colours, breaking the rule
that screens never hardcode colour. Justify the exception.
\ans{Google's brand terms forbid recolouring the ``G''. Routing it through the
palette would recolour it on every theme and put the app out of compliance, so
the literals are the correct behaviour and are commented as such to stop a future
contributor ``fixing'' them. The GitHub mark is monochrome and carries no such
restriction, so it takes a tint that defaults to the text token and inverts
properly on the light themes.}

\item Dynamic Type is supported through semantic text styles, with exactly one
fixed-size exception. Where is it, and why was that exception acceptable?
\ans{The exception is the numeral inside the score ring, which uses a fixed-size
font. The ring has a fixed diameter, so a numeral that scaled with the user's
text size would overflow it at the larger accessibility settings. Everything that
carries meaning as prose uses a semantic style and scales; the alternative would
be making the ring itself scale, which is the better fix and simply was not
worth it for one component.}

\end{questions}

\section{Backend and API Design}
\begin{questions}[resume]

\item Walk through the layering of the backend: routes, services, schemas and
models. What belongs in each, and what rule stops logic leaking between them?
\ans{Routes own HTTP concerns only --- the auth dependency, status codes, and
turning a service result into a response. Services own the actual work and the
provider calls. Schemas are the Pydantic request and response boundary. The
models are the SQLModel tables, which live in the separate database package. The
rule that holds it together is that a route never talks to a language model
directly and a service never sees a request object, so each layer is replaceable
without touching its neighbours.}

\item Why is Pydantic used as a boundary layer when SQLModel can already serialise
table models directly?
\ans{Because a table model exposes every column, including the password hash and
the email verification code. Returning one directly means a new column is
published to clients the moment it is added, which is a data leak waiting on a
future commit. An explicit response schema states exactly what leaves the server,
and it also decouples the wire format from storage, so a column can be renamed
without breaking every installed app.}

\item Authentication is a dependency injected into every protected route rather
than middleware. What are the advantages of that approach in FastAPI?
\ans{It is opt-in per route, so public endpoints simply do not declare it and
there is no allowlist of paths to keep in step. It appears in the generated
OpenAPI schema, so the docs show which endpoints are protected. It returns a
typed value --- the user identifier --- straight into the handler signature rather
than stashing it on the request. And it can be replaced wholesale in tests
through FastAPI's dependency override mechanism.}

\item Some endpoints, such as the server-side question bank, are deliberate
unimplemented stubs. Why were they left in the codebase rather than deleted?
\ans{They record a decision. An endpoint that raises \code{NotImplementedError}
with a comment saying the quizzes are bundled client-side makes it clear the
option was considered and rejected, whereas an absence looks like an oversight
and invites someone to build it. The cost is dead code, which is mitigated by
listing them explicitly in the project documentation. If they were still unbuilt
in a year, deleting them would be the right call.}

\item Social sign-in is brokered entirely by the backend, so no client identifiers
or secrets are compiled into the app. Walk through the flow and explain what
attack this prevents.
\ans{The app opens the backend's login endpoint in an authentication session; the
backend redirects to the provider; the provider calls back to the backend, which
exchanges the code for a token using its own secret, finds or creates the user,
issues our JWT, and redirects to the app's custom scheme, which the session
intercepts. The attack it prevents is trivial secret extraction: an app binary is
readable by anyone who downloads it, so a client secret shipped inside one is a
published secret, and an attacker could impersonate the app to the provider.}

\item The cross-site request forgery \code{state} parameter in the OAuth flow is a
signed JSON Web Token rather than a random string in a session. What does that
buy you, and what does it cost?
\ans{It makes the check stateless. The server can verify the value it gets back
using only its signing key, so nothing has to be stored between the two legs of
the flow --- which means a restart mid-login does not break it and it would work
unchanged across several instances. The cost is that it cannot be revoked before
it expires, and it is larger on the wire than an opaque random string.}

\end{questions}

\section{Data Model and Persistence}
\begin{questions}[resume]

\item There is no migrations framework. New columns are added by a hand-written
\code{ALTER TABLE} routine that runs on startup. Defend that decision, and state
the point at which it stops being tenable.
\ans{For a single instance with a handful of tables, a migration tool is real
setup plus ceremony on every change, and the routine here is genuinely safe: it
adds only nullable columns, checks \code{PRAGMA table\_info} before each one so it
is idempotent, and leaves existing rows to read as NULL. It stops being tenable
at the first change that is not an additive nullable column --- a rename, a type
change, a new NOT NULL column, or anything needing a data backfill --- and
immediately if there is ever more than one process writing, because the startup
routine would race.}

\item Interview transcripts are stored as a serialised JSON string rather than in
a related turns table. What are you trading away, and what query would force you
to normalise?
\ans{Querying. As it stands you cannot ask the database anything about the inside
of a conversation --- every answer requires loading whole rows and unpacking them
in Python, and nothing about a turn can be indexed. What is gained is one row per
interview, an atomic append, and no join on the hot path. Any analytics question
would force the change: average answer length by round type, or how far students
typically get before abandoning.}

\item Identity throughout the app is by stable identifier rather than by array
position --- option letters, quiz identifiers, LeetCode slugs. What concrete bug
class does that prevent?
\ans{Silent corruption of saved progress when content is edited. If solved state
were keyed by row number, inserting one problem into a company list would shift
every mark below it and quietly tell the user they had solved things they had
not. Because the key is the LeetCode slug, rows can be reordered freely, and
because it is the slug rather than a per-file identifier, solving a problem marks
it solved in all of the company lists that contain it.}

\item Quiz difficulty is a property of the quiz, not of individual questions, and
the validator rejects a per-question difficulty. Why was that constraint worth
enforcing in a script?
\ans{Because JSON is not type-checked, so a stray per-question difficulty field
would be perfectly valid JSON, would be silently ignored by the Swift decoder,
and would leave an author believing they had set something that had no effect.
Keeping difficulty in one place is what stops a quiz labelled Hard from showing an
Easy badge on question three. The validator also cross-checks category and subject
values against the actual Swift enumerations, which catches the worse case: a
value that parses as JSON but fails to decode at runtime.}

\item Progress stores are write-through caches: they update local storage
immediately and mirror to the server afterwards. How do you handle a device that
was offline for the whole session, and how do conflicts resolve?
\ans{The local write is the source of truth for the session, so an offline
student sees every change instantly and loses nothing. On the next successful
sign-in the sync routine pulls server state and pushes anything that exists only
locally, which also serves as the migration path for progress made before the
backend existed. Conflicts resolve by merge rather than by overwrite: solved
problems take the union of both sides, and quiz scores keep the higher one, so
neither device can erase the other's work.}

\item Each local store carries an owner user identifier. What happens when a
second student signs in on the same device, and why was that not left to the
operating system?
\ans{On sign-in the store compares the stored owner with the incoming user and
discards the cache outright if they differ, rather than adopting it. Without that
check the second student would inherit the first one's solved problems and quiz
scores, which is both a privacy leak and visibly wrong. It cannot be delegated
because iPhone has no multi-user account model --- as far as the operating system
is concerned there is one user of the device, so the application has to draw the
boundary itself.}

\end{questions}

\section{Artificial Intelligence and Language Model Integration}
\begin{questions}[resume]

\item The system uses Groq as the primary provider with Gemini as a fallback,
behind a single interface. Why two providers, and how does the calling code stay
unaware of which one served a request?
\ans{Two providers because a single free tier is not a dependable base for the
one feature the project is judged on, and because Groq is fast and also serves
the speech-to-text model. Everything funnels through one internal generate
helper, so routes call functions like \code{generate\_follow\_up} and never see a
provider name, a model identifier or an HTTP client. Reordering or replacing a
provider is a change to the configured chain and the per-provider call functions,
and nothing above that layer moves.}

\item On a rate-limit response the service walks a chain of models rather than
failing. Explain the quota structure that makes this effective rather than merely
a retry.
\ans{The free tier's daily cap is enforced per model, per project --- the quota
identifier is literally scoped that way --- so an exhausted model says nothing
about the next one, which has its own untouched allowance. That is what makes the
walk genuinely additional capacity rather than a retry against the same empty
bucket. It is also why the budget is worth stating plainly: one interview costs a
call to start plus one per answer, and a resume summary costs one more.}

\item Provider errors are never passed through verbatim to the user. What does the
failure path look like instead, and why does that matter for this audience in
particular?
\ans{Failures are mapped to a small set of clean statuses --- a transport or model
error becomes a 502, a missing key or fully exhausted quota becomes a 503 --- each
carrying one short human sentence. It matters because the raw quota response is a
forty-line diagnostic dump, and the person receiving it is a nervous student in
the middle of a mock interview who can do nothing with it. An error message is
part of the product, not a debugging artefact.}

\item Interview replies signal termination with a sentinel marker, while the
end-of-round assessment is requested as JSON. Justify using two different
mechanisms in the same system.
\ans{Because the two outputs have opposite constraints. An interview reply is
spoken aloud, so the prompt forbids markdown, numbering and punctuation
furniture; demanding JSON in the same breath would fight those rules and produce
worse speech. A sentinel is one token appended to the end and stripped before the
text is returned, so the spoken output stays clean. The debrief is never spoken
--- it is rendered as a report with four typed fields --- so JSON is free there and
strictly better than parsing prose.}

\item The interview can end itself when a candidate answers in bad faith. How do
you distinguish bad faith from a nervous, wrong or garbled answer, and which way
does the system err?
\ans{The prompt states both cases explicitly, and the list of things that must
\emph{not} end an interview is deliberately longer than the list that must:
wrong, blank, rambling, one-word, memorised, ``I don't know'', and mangled
speech-to-text all keep it running. Ending is reserved for threats, abuse,
discriminatory content, plain trolling, or an attempt to instruct the model to
change its behaviour. It errs firmly toward continuing, because every user is a
student who is bad at interviews, and a tool that walks out on a nervous answer is
worse than no tool.}

\item A candidate could attempt prompt injection to obtain a favourable result.
What defences exist, and what are their limits?
\ans{Instructing the model to award a result is treated as bad faith and ends the
round, and in testing an injection attempt drew a warning and then a termination
when repeated. The limit is that this is a prompt-level defence, so a sufficiently
novel phrasing may get through --- it should be described as mitigation, not a
guarantee. The structural defences are stronger: the score is produced by a
separate call over a stored transcript rather than by the interviewer in
conversation, and termination is recorded in the database, so a model that is
talked around in one turn still cannot revive a closed round.}

\item The round is closed on the server, not only in the interface, and further
answers are rejected. Why was interface-level termination insufficient?
\ans{Because client-side state is advisory. Ending the round only in the app means
the session row is still marked in progress, so a replayed request --- from a
retry, a stale view, or anyone with the token and a HTTP client --- resumes an
interview the interviewer had walked out of. Setting the status server-side and
returning a conflict on any further answer makes the ending a fact about the
system rather than a fact about one screen.}

\item The interview is capped at a fixed number of rounds by the client rather
than the model deciding when enough ground is covered. Explain the reasoning.
\ans{The model has no idea what the product considers a complete interview, and
left to choose it tends to wrap up after two or three questions, which is not
practice. The cap is a product decision about how long a useful session is, so it
belongs where product decisions live. The one place a natural ending is genuinely
available is the group discussion, whose moderator has an in-fiction reason to
close the floor --- so that round, and only that round, is allowed to end itself.}

\item In the interview prompts the target role is placed above the resume material
and stated as outranking it. What behaviour were you correcting, and how did you
verify the change worked?
\ans{Originally the role appeared once, as a noun in the opening sentence, which
left the model free to interview a data-science candidate about generic web
backend trivia --- the role was present but carried no weight. The fix was a
dedicated block above each round's brief instructing the model to work out what
the role does day to day and ask only questions serving it. It was verified by
generating the same fundamentals round for two roles: a data scientist was asked
about hash tables versus balanced trees over large datasets, and an embedded
engineer about a process blocking on a held resource.}

\end{questions}

\section{Voice, Resume and On-Device Processing}
\begin{questions}[resume]

\item Apple's on-device speech recogniser was removed in favour of uploading audio
for transcription. What forced that, and what did the app lose as a result?
\ans{It cannot initialise in the Simulator --- no on-device model ships there, and
the server-based recogniser fails outright --- which made the entire voice feature
untestable without physical hardware on every change. What was lost is the live
word-by-word transcript: text now exists only after the recording is uploaded and
returned, so the interface shows a recording state rather than words appearing.
Reinstating it would mean gating on the on-device support flag and keeping the
upload path as the fallback.}

\item Resume text extraction prefers the PDF's embedded text layer and falls back
to Vision optical character recognition only when there is none. Why is that
ordering correct?
\ans{An embedded text layer is the document's own characters, so reading it is
lossless; running recognition over a rendering of the same page can only
introduce misreads. Nearly every resume is exported from Word, Docs or LaTeX and
carries one, so the fast, exact path is also the common path. Recognition earns
its place only on scanned or image-only resumes, where there is no text layer to
prefer --- and the code decides by measuring what it actually extracted rather
than by trusting the file type.}

\item A real resume was mis-parsed because the PDF text layer glued a section
heading onto the end of the previous line. How was that diagnosed, and why did
optical character recognition not solve it?
\ans{By extracting the document and printing which lines the parser classified as
headings: the projects heading never appeared, because the text layer emitted
``PROJECTS \& PUBLICATIONS'' welded to the tail of the preceding bullet, so the
whole-line heading match could never fire. Recognition did not help --- it read
the heading and the line above as one visual row and garbled the overlap into
nonsense, which is worse than the original. The fix belonged in the parser:
split a trailing capitalised heading back onto its own line, trying the longest
known heading first so the combined heading is not mistaken for the shorter one
inside it.}

\item Only the projects section of a resume is transmitted, after contact details
are stripped. What threat model motivated that, and where exactly does the
redaction happen?
\ans{The free tier's terms permit submitted content to be retained and reviewed by
humans, and a resume is the densest personal data a student owns --- name, email,
phone, address, education history. The redaction happens entirely on the device,
before any request is made: the parser narrows the text to the projects section
and strips email and phone patterns from what remains. Nothing resume-derived is
stored server-side at all; the summary endpoint is deliberately stateless and the
result is kept only in local storage on the phone.}

\item The company problem lists are parsed with a hand-written comma-separated
value reader rather than by splitting on commas. What in the data made the naive
approach incorrect?
\ans{The topics column is a quoted list that itself contains commas --- a row ends
with \code{"Array, Hash Table"} --- so splitting on every comma produces an extra
field and shifts the parse for that row. Since topics is the last column, the
damage is a corrupted topic list on precisely the rows with more than one topic,
which is most of them. The reader implements RFC~4180 quoting, which is a short
state machine and the only correct way to read the format.}

\end{questions}

\section{Security and Privacy}
\begin{questions}[resume]

\item Describe the full account lifecycle: registration, one-time-password
verification, token issue, session validation and sign-out. Where is the token
stored on the device, and why there?
\ans{Sign-up creates the user and issues a six-digit code by email, which a verify
endpoint confirms; until then the app holds the student on a verification screen.
Login returns a JWT, which is stored in the Keychain rather than in user defaults
--- defaults are a plain property list, readable from a file-system backup,
whereas the Keychain is encrypted and access-controlled. On each launch the stored
token is validated against the profile endpoint before the app trusts it, and
sign-out clears both the token and the per-user caches.}

\item A token that expires mid-session must not strand the user. How is an
unauthorised response handled, and how does that differ from a network failure at
launch?
\ans{The network layer exposes an unauthorised callback that fires on a 401 for
any authenticated request and returns the app to the login screen, so an expired
token surfaces as a sign-in prompt rather than as a broken screen. At launch the
distinction is deliberate and important: a genuine 401 clears the stored token,
but a network error keeps it, so a student opening the app on a train is not
signed out for being offline. Treating those two identically is the common bug,
and it is the one this design exists to avoid.}

\item The application transport security exception permitting cleartext local
networking is present for development. What is the risk if it reaches a release
build, and how is that prevented?
\ans{The exception used is the narrow one --- it permits cleartext only to local
network hosts and cannot enable arbitrary plaintext traffic to the internet, so
even shipped it does not expose production traffic. It is needed because the
debug build points at a local server over HTTP. The structural protection is that
the base URL is chosen by build configuration: debug points at localhost and
release points at the deployed HTTPS host, so no release traffic is cleartext
regardless of what the transport policy would permit.}

\end{questions}

\section{Testing, Deployment and Scale}
\begin{questions}[resume]

\item There is currently no automated test suite. Given that, how was correctness
established, and what would you write first if you had one week?
\ans{By compiling both halves on every change, driving the API end to end against
a running server with the model call stubbed so results were deterministic, and
using live calls only for behaviour that nothing but a real model can demonstrate
--- that a hostile answer ends a round while a nervous one does not, for instance.
Given a week I would start where the risk is highest and the tests are cheapest:
route tests with the model patched, covering authentication, the caching of the
debrief, the conflict on an already-ended session and the sentinel parsing; then
pure unit tests for the resume parser and the CSV reader against real fixture
files, since both are parsers with known edge cases and no dependencies. The quiz
content validator already acts as a gate and would move into continuous
integration on day one.}

\item The service runs on a single instance with SQLite on a mounted volume.
Identify the first bottleneck under load and describe your migration path.
\ans{The first bottleneck is that SQLite has a single writer and the volume pins
the application to one machine, so it cannot be scaled horizontally at all --- the
write lock, not CPU, is what gives way first. The path is to move to Postgres,
which is mostly a connection-string change because the tables are defined through
an ORM, plus adopting a real migration tool at the same time, since the
hand-rolled column routine does not survive multiple writers. Once storage is
external the application layer is stateless and scales out normally. The second
bottleneck is the language-model quota, which is scoped per project and therefore
completely unaffected by adding instances; that needs a paid tier or a request
queue with backpressure, and it is worth naming because it is the one that
scaling the servers cannot fix.}

\end{questions}

\vfill
\noindent\color{muted}\small
Prepared as revision material. Every answer is verifiable against the project's
source and its accompanying documentation.
\color{ink}

\end{document}
